\section{The theorem and the precise claim}
\label{sec:statement}

The MathWorld formulation concerns expressions generated by rational constants,
$\pi$, $\log 2$, the variable $x$, addition, multiplication, composition, sine,
exponential, and absolute value. It says that no algorithm decides whether the
represented function is zero for every real input \cite{mathworld}. Negative
rational constants make subtraction available. All these expressions represent
total continuous real functions; there is no domain-of-definition issue here.

Let $\mathcal R$ denote this class. Composition can be eliminated from the
syntax by substituting the inner expression for the variable, so we use the
equivalent grammar
\[
 E ::= q\mid x\mid\pi\mid\log 2\mid E+E\mid E\cdot E
       \mid\sin E\mid\exp E\mid |E|,
 \qquad q\in\Q.
\]
Let $\mathcal S$ be the sublanguage obtained by omitting $\log 2$ and $\exp$.
Our constructed expressions have the still more special form
\begin{equation}\label{eq:restricted-form}
 E_P(t)=|A_P(t)|-A_P(t),
 \qquad
 A_P\in\Q[\pi,t,\text{iterated sine expressions}],
\end{equation}
where $A_P$ contains no absolute value. In particular there is exactly one
absolute-value occurrence in the displayed expression tree. We do not claim
undecidability of identity after absolute value is removed.

\begin{theorem}[Richardson in arithmetic]\label{thm:main}
There is an explicitly specified primitive recursive rational approximation
function $\App$ for expression codes. With $\Id$ and $\Dec$ defined in
Sections~\ref{sec:arithmetic} and \ref{sec:diagonal}, respectively,
\[
 \ipra\ \vdash\ \forall d\,\neg\Dec(d).
\]
Here $\Dec(d)$ says that the ordinary program with code $d$ terminates with a
Boolean answer on every valid expression and answers correctly for identity
to zero. It quantifies over arbitrary program codes; it does not restrict the
claim to algorithms whose termination is already provable in iPRA.

The same nonexistence statement holds when correctness and termination are
required only on the family \eqref{eq:restricted-form}. In the standard
interpretation of arithmetic, $\Id(e)$ is equivalent to the usual assertion
that the real function represented by $e$ vanishes everywhere.
\end{theorem}

The theorem is an assertion about what iPRA proves, not an assertion that
iPRA contains a set of all real numbers. Its first-order sentence quantifies
only over natural-number codes: expressions, rational inputs, precision
parameters, and finite computations. The last sentence of the theorem
identifies the ordinary mathematical meaning of those codes.

\subsection{The central reduction}

The construction below proves a stronger, useful intermediate statement.
For every integer polynomial $P$ in $n\ge1$ variables, one can compute an
expression $E_P\in\mathcal S$ such that
\begin{equation}\label{eq:poly-negative-summary}
 \ipra\ \vdash\quad
 \Id(E_P)\ \longleftrightarrow\
 \neg\exists \boldsymbol m\in\N^n\;P(\boldsymbol m)=0.
\end{equation}
The notation is schematic: polynomials and tuples are encoded by natural
numbers, and the uniform version is a single first-order sentence. There
are primitive recursive maps from a root to a strict rational nonzero
certificate and back. The analytic proof therefore provides positive finite
witnesses before taking negations.

For one fixed universal computation predicate, formalized MRDP gives a fixed
polynomial $P_0(e,\boldsymbol m)$. Its negative translation yields
\begin{equation}\label{eq:negative-halt-summary}
 \ipra\ \vdash\quad
 \neg\Halt(e)\ \longleftrightarrow\
 \neg\exists\boldsymbol m\;P_0(e,\boldsymbol m)=0.
\end{equation}
Specializing $P_0$ at $e$ and applying \eqref{eq:poly-negative-summary} gives
a primitive recursive expression constructor $C$ with
\begin{equation}\label{eq:main-reduction}
 \ipra\ \vdash\quad
 \Id(C(e))\ \longleftrightarrow\ \neg\Halt(e).
\end{equation}
A direct two-branch diagonal argument then proves Theorem~\ref{thm:main}.

\subsection{Dependency and verification boundaries}

The deep imported theorem is the \emph{formalized} MRDP theorem of
Gaifman--Dimitracopoulos \cite{gd}, in the precise instance stated in
Section~\ref{sec:foundation}. A standard-model MRDP theorem alone would not
establish \eqref{eq:negative-halt-summary}. The weak-arithmetic formalization
is independently recorded in the primary research literature
\cite{cmy} and in \cite{jerabek}.

\begin{center}
\begin{tabularx}{\linewidth}{@{}p{0.30\linewidth}X@{}}
\toprule
Part & Status in this deliverable\\
\midrule
Polynomial-to-expression map & Explicit formula and exact-arithmetic reference implementation.\\
Analytic gap and rational density & Complete quantitative human proofs, with primitive recursive bounds.\\
Negative MRDP bridge & Proved by translating the cited established EFA theorem; no new MRDP proof is claimed.\\
Diagonal argument & Complete constructive proof and fully written parametric Lean logical core.\\
Final internal Lean theorem & Detailed definitions and implementation plan; not completed or kernel-checked.\\
\bottomrule
\end{tabularx}
\end{center}

The distinction in the last two rows is deliberate. A Lean theorem over
ordinary natural numbers, or a Lean theorem parameterized by representation
lemmas, is not itself a derivation in the object theory iPRA. The requested
human proof can use established mathematical theorems; a completed machine
formalization must additionally supply their proof objects and the indicated
arithmetical translations.
